\part{Interpretation Before the Sovereign Subject}

\chapter{Reading \emph{The Will to Power} as a Problematic Archive}

The textual basis of any account of Nietzschean interpretation must first settle what is being read. \emph{The Will to Power} appeared as a book under Nietzsche's name only after his collapse, assembled by his associates and his sister from notebook material of the 1880s, and the numbered aphorisms of the widely used English edition \citep{kaufmann1967} follow an arrangement that Nietzsche did not give them. The critical edition prepared by Colli and Montinari restored the notebooks to their chronological order and made it possible to see that the plan of a systematic work under this title was entertained, revised and eventually abandoned in 1888, when Nietzsche turned instead to the writings he completed \citep{montinari1982}. A passage quoted from the compilation therefore comes with two kinds of uncertainty. It may record a thought that was later developed, tested against objections and published in a finished work, or it may record a trial formulation that Nietzsche dropped.

The method adopted here follows from that situation. Passages from the notebooks are used to reconstruct a constellation of problems and are never treated as articles of a finished doctrine. Where a claim about interpretation appears in the compilation, it is located in its notebook context where that can be done, and compared with its treatment in \emph{The Gay Science}, \emph{Beyond Good and Evil}, \emph{On the Genealogy of Morality} and \emph{Thus Spoke Zarathustra} \citep{nietzsche1974,nietzsche1966,nietzsche1967gm,nietzsche2006}. Agreement between the notebooks and a published work strengthens an attribution. Disagreement is recorded as a tension or as a change of emphasis, and is not smoothed away to produce a unified position. Systematic readings of Nietzsche exist and are argued with care \citep{richardson1996}, and the present method does not exclude them. It requires only that any systematic reading say which of its premises rest on published statements and which on material Nietzsche chose not to publish.

The comparison can be made concrete for the doctrine of the next chapter. The notebook claim that knowing is a schematizing that imposes regularity as far as practical needs require (section 515 of the compilation) has published counterparts of different weight. Section 24 of \emph{Beyond Good and Evil} describes cognition as a will to simplification that has contrived to preserve ignorance in the service of life, and section 354 of \emph{The Gay Science} derives conscious thought from the pressure of communication, so that what reaches awareness is already shaped by what can be shared. Both passages support the notebook's picture of cognition as practical simplification, and neither states it as a general theory of knowledge in the compilation's terms. The attribution adopted below is accordingly that Nietzsche held the schematizing thesis as a recurring problem and not that he fixed it as a doctrine, and the dating and notebook identifiers of the relevant entries in the critical edition remain to be supplied wherever a claim rests on the compilation alone.

A literary question enters at the outset and recurs throughout the book. Aphoristic and experimental writing, when reorganized into the form of a philosophical system, acquires expectations of continuity, hierarchy and doctrinal consistency that the original writing did not claim to satisfy. A reader who meets numbered sections arranged under headings on knowledge, nature, society and art will tend to assume that earlier sections ground later ones and that apparent contradictions are failures of exposition. The editorial arrangement thus supplies an interpretive frame before any sentence is read. This is itself an instance of the phenomenon under study, since an organization imposed on material determines which relations among its parts are visible, which are salient and which are treated as noise. The material form of Nietzsche's reception is accordingly part of the argument about how interpretation organizes its objects, and the reading practice adopted in later chapters aims to resist the frame without pretending to a neutral standpoint outside any frame.

\chapter{Knowledge as Schematization}

Nietzsche's treatments of identity, regularity and the concept share a recurring claim: that the cognition which makes a stable world available does so by treating unlike things as like and by imposing regularity on a flux that does not present itself in that form. The notebooks put this in terms of schematization, an activity that, as section 515 of the compilation has it, imposes upon chaos as much regularity and form as practical needs require. The question this chapter pursues is whether such activity should be described as the discovery of an order or as the making of experience tractable through an ordering. The two descriptions differ in what they permit one to say about error. If cognition discovers order, a schema is false to the degree that it fails to match what is there. If cognition makes experience tractable, a schema may be adequate or inadequate relative to a task, a scale and a form of life, and the word error needs a different home.

The question of error is sharper than it first appears. A useful simplification does not become false merely because it suppresses differences, since every concept suppresses some. What matters is whether the suppressed differences are idle with respect to the purposes the concept serves. The same simplification can be adequate for one purpose and ruinous for another, and it can pass from the first condition to the second when circumstances change. A reading of Nietzsche on which all schematization is equally falsifying cannot make this distinction, and a reading on which it is simply true cannot explain why he treats the same schemata as both indispensable and as objects of suspicion.

Coarse-graining supplies a precise vocabulary for the point. Let $X$ be a detailed state space and let an interpretive organization be a map
\[
\phi : X \longrightarrow Z,
\]
where $Z$ contains the distinctions available to the organism. Because distinct states of $X$ may be sent to the same state of $Z$, the organization necessarily discards differences. The relevant question is which discarded differences remain irrelevant to action and which later become consequential. This can be made exact relative to a task.

\begin{definition}[Task-relative sufficiency]
Let $T$ be a task specified by a class of admissible policies and a loss $\ell$ defined on $X$. An organization $\phi:X\to Z$ is \emph{$\varepsilon$-sufficient for $T$} if the best expected loss attainable by policies that depend on $X$ only through $\phi$ exceeds the best expected loss attainable by policies that depend on $X$ directly by no more than $\varepsilon$.
\end{definition}

The definition gives sense to the claim that an organization is adequate without being a copy. It also makes adequacy fragile in the way the notebooks suggest. If the task changes, or if the distribution of encountered states shifts, a map that was $\varepsilon$-sufficient for the earlier task and distribution may cease to be so, although nothing about the map itself has altered. A simplification can thus pass from harmless to harmful through a change in what the organism does or where it lives. This gives a precise reading of the claim that errors of cognition are conditions of life: the simplification is a condition of life for the tasks it serves, and an error for those it does not.

\begin{remark}
The definition does not say that all schematization is innocent. It says that the evaluation of a schema requires a specification of task and environment, and that the specification can itself be contested. A reconstruction of Nietzsche's view that left this contestation out would explain only the half of his position on which schemata are useful.
\end{remark}

\chapter{The Ego as an Interpretation of Activity}

Nietzsche's criticism of a substantial subject, a simple inner proprietor who thinks, wills and owns its acts, is often taken to imply a denial of agency. The two should be kept apart. The disappearance of a simple proprietor leaves the fact of organized action in place, together with the need to explain how a plurality of processes acquires enough coordination to act, to remember, to promise and to experience itself as continuous. The explanatory task is thereby changed rather than abolished. Instead of asking what the self is, one asks how a self-describing coordination is produced and maintained, and what its description leaves out. Nietzsche's own positive proposal, in \emph{Beyond Good and Evil}, section 12, is to understand the soul as a multiplicity, a social structure of drives and affects, and section 19 analyses willing as a relation of commanding and obeying within the one who wills. The feeling of being a unified agent can conceal the complex activity that produces it, since the one who commands identifies with the commanding side and takes credit for what the whole organization accomplishes. The ego is accordingly partly an achievement of organization and partly a misleading interpretation of that achievement.

The comparison with contemporary modelling concerns the distinction between a process and its representation of its own activity. A system may maintain a model that summarizes its bodily states, its action possibilities and the expected consequences of acting, without that model transparently disclosing the mechanisms that generate the behaviour it summarizes. Such a self-model is useful in the way schemata are useful: it compresses a large and partly inaccessible set of processes into a form that can guide planning and be communicated to others. It can also mislead in a systematic way. A model that presents the system as a unified deliberator will explain actions by appeal to decisions and will render the contest of concerns that produced them invisible, which is the pattern that Nietzsche diagnoses in the ordinary grammar of willing.

The easy conclusion, that an interpreted self is therefore unreal, does not follow and should be resisted. An interpretation can have causal consequences. It can regulate conduct, structure what the organism attends to, shape what others expect of it, and become an object of investment, so that challenges to it are felt as challenges to the organism and not merely to a proposition. These effects obtain even when the interpretation's account of its own origins is mistaken. The self so understood is neither a substance nor an illusion in the sense of a mere appearance without effect. It is an organization that has consequences, and the questions it raises are about its authority, its revisability and the cost of revising it, which later chapters take up.

The distinction between an agent and the model an agent has of itself also bears on the relation between responsibility and metaphysics. Promising and being held to a promise require that conduct be coordinated across time, and a self-model that supports such coordination is practically necessary whatever its metaphysical standing. The Nietzschean question is therefore what kind of coordination a given self-interpretation makes possible and what kinds it forecloses, and this question can be asked without settling whether any inner proprietor exists.

\chapter{Drives, Perspectives, and Competing Interpretations}

The term drive is used by Nietzsche in ways that resist a single analysis. Sometimes it names an impulse that presses for discharge, sometimes an evaluative orientation that gives things their significance, sometimes a competing organization that seeks to interpret the whole in its own terms. These senses are not obviously consistent, and the comparative problem is sharpened by a difference of method. Many active-inference models begin with a specified agent and a specified objective, and ask how the agent should perceive and act. Nietzsche frequently makes the organization of the agent itself a site of conflict, so that the question of whose objective governs is open and its answer can change.

A reconstruction can represent partially competing evaluative processes through objectives $F_i$, each of which assesses states and policies from its own standpoint, and an effective organization such as
\[
F_{\mathrm{eff}}(s) = \sum_i w_i(s)\,F_i(s).
\]

\begin{proposal}
Treat a Nietzschean multiplicity as a family of objectives $\{F_i\}$ together with state-dependent weights $w_i(s)$ that determine which concerns govern conduct in which situations.
\end{proposal}

What matters philosophically is the weights and not the sum. The sum represents a momentary settlement, whereas the weights encode a hierarchy, and the question that Nietzsche presses is what determines which concerns dominate and how that determination can change. If the weights are fixed, the multiplicity collapses into a single objective with a complicated shape, and nothing distinctively Nietzschean remains. If the weights depend on state, they can shift as circumstances shift, which captures the idea that a drive may dominate in one setting and be subordinate in another. If they depend on the history of the organism, a further structure is needed to describe how hierarchies are acquired, and this is where an account of formation and revision becomes necessary.

It is important to say what this proposal does not claim. It does not translate drives into computational modules, and it does not suggest that Nietzsche had a weighted sum in mind. It offers a form that makes certain commitments inspectable, and one of its uses is to reveal where the form fails. When the objectives $F_i$ are not commensurable, so that no common scale licenses adding them, the scalar $F_{\mathrm{eff}}$ is an artefact of the reconstruction and not a feature of the phenomenon. In such cases the conflict is better represented as an ordering that is partial, or as a contest whose resolution is a matter of the organization taking one form rather than another. Part of the later discussion of value creation turns on whether a framework that begins from a scalar objective can represent such conflicts without silently reintroducing an evaluator that stands above them.
